\documentclass[10pt,a4paper]{amsart}
\usepackage{amsmath,amssymb,amsthm,graphicx}
% Retained minimal preamble for the Stacks Project chapter more-algebra.tex (More on Algebra).
% The chapter file pinned-original-source.tex carries only \input{preamble} on its first line, so no
% preamble text lives inside the chapter itself.  The Stacks Project uses one single shared preamble for
% every chapter; the snapshot stages it identically as sources/stacks-chow/preamble.tex and
% sources/stacks-etale-cohomology/preamble.tex.  The block below is copied mechanically, line for line,
% from sources/stacks-chow/preamble.tex lines 14-18 (the Xy-pic setup) and 154-217 (the native theorem
% declarations together with the mathematical macros).  Nothing was added, removed or redefined; the
% wrapper class and the amsmath/amssymb/amsthm/graphicx base packages are supplied by the assembler.
% Chapter: more-algebra.tex -- More on Algebra.
% For commutative diagrams we use Xy-pic
\usepackage[all]{xy}

% We use 2cell for 2-commutative diagrams.
\xyoption{2cell}

% Theorem environments.
%
\theoremstyle{plain}
\newtheorem{theorem}[subsection]{Theorem}
\newtheorem{proposition}[subsection]{Proposition}
\newtheorem{lemma}[subsection]{Lemma}

\theoremstyle{definition}
\newtheorem{definition}[subsection]{Definition}
\newtheorem{example}[subsection]{Example}
\newtheorem{exercise}[subsection]{Exercise}
\newtheorem{situation}[subsection]{Situation}

\theoremstyle{remark}
\newtheorem{remark}[subsection]{Remark}
\newtheorem{remarks}[subsection]{Remarks}

\numberwithin{equation}{subsection}

% Macros
%
\def\lim{\mathop{\mathrm{lim}}\nolimits}
\def\colim{\mathop{\mathrm{colim}}\nolimits}
\def\Spec{\mathop{\mathrm{Spec}}}
\def\Hom{\mathop{\mathrm{Hom}}\nolimits}
\def\Ext{\mathop{\mathrm{Ext}}\nolimits}
\def\SheafHom{\mathop{\mathcal{H}\!\mathit{om}}\nolimits}
\def\SheafExt{\mathop{\mathcal{E}\!\mathit{xt}}\nolimits}
\def\Sch{\mathit{Sch}}
\def\Mor{\mathop{\mathrm{Mor}}\nolimits}
\def\Ob{\mathop{\mathrm{Ob}}\nolimits}
\def\Sh{\mathop{\mathit{Sh}}\nolimits}
\def\NL{\mathop{N\!L}\nolimits}
\def\CH{\mathop{\mathrm{CH}}\nolimits}
\def\proetale{{pro\text{-}\acute{e}tale}}
\def\etale{{\acute{e}tale}}
\def\QCoh{\mathit{QCoh}}
\def\Ker{\mathop{\mathrm{Ker}}}
\def\Im{\mathop{\mathrm{Im}}}
\def\Coker{\mathop{\mathrm{Coker}}}
\def\Coim{\mathop{\mathrm{Coim}}}

% Boxtimes
%
\DeclareMathSymbol{\boxtimes}{\mathbin}{AMSa}{"02}

%
% Macros for moduli stacks/spaces
%
\def\QCohstack{\mathcal{QC}\!\mathit{oh}}
\def\Cohstack{\mathcal{C}\!\mathit{oh}}
\def\Spacesstack{\mathcal{S}\!\mathit{paces}}
\def\Quotfunctor{\mathrm{Quot}}
\def\Hilbfunctor{\mathrm{Hilb}}
\def\Curvesstack{\mathcal{C}\!\mathit{urves}}
\def\Polarizedstack{\mathcal{P}\!\mathit{olarized}}
\def\Complexesstack{\mathcal{C}\!\mathit{omplexes}}
% \Pic is the operator that assigns to X its picard group, usage \Pic(X)
% \Picardstack_{X/B} denotes the Picard stack of X over B
% \Picardfunctor_{X/B} denotes the Picard functor of X over B
\def\Pic{\mathop{\mathrm{Pic}}\nolimits}
\def\Picardstack{\mathcal{P}\!\mathit{ic}}
\def\Picardfunctor{\mathrm{Pic}}
\def\Deformationcategory{\mathcal{D}\!\mathit{ef}}

% --- appended 2026-10-12 (Chief): the project own definitions of the reference/slogan/history
% display environments, copied mechanically from the shared Stacks preamble lines 8-12
% (\usepackage{verbatim} plus the three \newenvironment lines, which the chapter preamble keeps
% verbatim from the project preamble). Nothing added, removed or redefined: these environments
% carry no mathematical content and the project itself typesets them as comments. Without these
% three lines a counted statement whose own text uses them fails with Environment slogan undefined.
\usepackage{verbatim}
\newenvironment{reference}{\comment}{\endcomment}
\newenvironment{slogan}{\comment}{\endcomment}
\newenvironment{history}{\comment}{\endcomment}
\usepackage{cleveref}
\crefname{originallemma}{original lemma}{original lemmas}
\Crefname{originallemma}{Original lemma}{Original lemmas}
\makeatletter
\newlabel{lemma-tensoring-Deligne-system}{{0G3N}{source}{}{sourcepointer-1}{}}
\newlabel{lemma-tensoring-Deligne-system@cref}{{[originallemma][0][]0G3N}{[1][0][]source}{}{}{}}
\makeatother
\crefname{originallemma}{original lemma}{original lemmas}
\Crefname{originallemma}{Original lemma}{Original lemmas}
\makeatletter
\newlabel{lemma-ext-annihilated}{{0927}{source}{}{sourcepointer-2}{}}
\newlabel{lemma-ext-annihilated@cref}{{[originallemma][0][]0927}{[1][0][]source}{}{}{}}
\makeatother
\title{A high power of an ideal maps through the derived tensor product of the ideal with the module}
\author{The Stacks Project Authors}
\date{stacks-project:more-algebra.tex:local-snapshot-20261012; GNU FDL 1.2 or later}
\begin{document}
\maketitle
\setcounter{subsection}{1}
\noindent\textit{Editorial note.} All delivered mathematical text is the author's own, copied from the pinned Stacks Project chapter file more-algebra.tex; no statement or proof step was rewritten, reordered or removed. The counted statement is the labeled lemma environment at source lines 29736-29742 (canonical tag 0928) and the counted proof is the proof environment at source lines 29744-29768, which immediately follows it and is its own proof.\par\medskip
\noindent\textit{Editorial note.} The printed text writes It can also been seen directly where It can also be seen is meant; the wording is copied exactly as printed and is not corrected.\par\medskip
\noindent\textit{Editorial note.} The preamble span runs to line 90, the current end of the retained preamble, which includes the appended environment definitions; the counted text does not use them, and the span is taken whole so that the file version is unambiguous.\par\medskip
\noindent\textit{Editorial note.} The source URL is the source-level Stacks Project URL, as required when more than one card is produced from one source; the per-claim tag is carried by claim\_key and by this note. The counted statement and proof cite no bibliography entry, so this card needs no bibliography.\par\medskip
\noindent\textit{Allowed standard prerequisites (editorial).}\begin{itemize}
\item For an ideal in a Noetherian ring the comparison between the derived tensor product with its powers and the ordinary products is a pro-isomorphism, which immediately gives a power factoring through the derived tensor product of the ideal itself. invoked at proof lines 29744-29745. Supplies a short proof of the statement, after which the author gives an independent direct argument.
\item The quotient by the ideal sits in a distinguished triangle with the ideal, so factoring through the derived tensor product with the ideal is equivalent to the vanishing of the composite into the derived tensor product with the quotient. used at proof lines 29746-29752. Reformulates the goal as the vanishing of one morphism in the derived category.
\item A finite module annihilated by an ideal can be absorbed by a sufficiently high power of that ideal in the Ext group, so an obstruction element can be killed after increasing the power. invoked at proof lines 29762-29765. Allows the successive factorization steps to be completed in finitely many rounds, which gives the required power.
\end{itemize}\medskip
\ifdefined\hypertarget\hypertarget{sourcepointer-1}{}\fi
\noindent\textit{Source reference.} Original lemma 0G3N: sources/stacks-more-algebra/pinned-original-source.tex, source lines 29677--29677. The reference is cited at proof line 29744 as the reason the statement holds, after which the author gives an independent direct argument; the direct argument, which is what the counted proof contains, uses this reference only as a remark and it is the subject of a separate card of this lane.\par\medskip
\ifdefined\hypertarget\hypertarget{sourcepointer-2}{}\fi
\noindent\textit{Source reference.} Original lemma 0927: sources/stacks-more-algebra/pinned-original-source.tex, source lines 29571--29571. The reference is applied at proof lines 29762-29765 to kill the obstruction element in the Ext group after increasing the power of the ideal; the induction that consumes it is inside the counted span.\par\medskip
\begin{lemma}
\label{lemma-factor-through-derived-tensor-product}
Let $A$ be a Noetherian ring. Let $I \subset A$ be an ideal. Let $M$
be a finite $A$-module. There exists an integer $n > 0$ such that
$I^nM \to M$ factors through the map $I \otimes_A^\mathbf{L} M \to M$
in $D(A)$.
\end{lemma}

\begin{proof}
This follows from Lemma \ref{lemma-tensoring-Deligne-system}.
It can also been seen directly as follows. Consider the distinguished triangle
$$
I \otimes_A^\mathbf{L} M \to M \to A/I \otimes_A^\mathbf{L} M \to
I \otimes_A^\mathbf{L} M[1]
$$
By the axioms of a triangulated category it suffices to prove that
$I^nM \to A/I \otimes_A^\mathbf{L} M$ is zero in $D(A)$ for some $n$.
Choose generators $f_1, \ldots, f_r$ of $I$ and let
$K = K_\bullet(A, f_1, \ldots, f_r)$ be the Koszul complex
and consider the factorization $A \to K \to A/I$ of the quotient map.
Then we see that it suffices to show that $I^nM \to K \otimes_A M$
is zero in $D(A)$ for some $n > 0$. Suppose that we have found an $n > 0$
such that $I^nM \to K \otimes_A M$ factors through
$\tau_{\geq t}(K \otimes_A M)$ in $D(A)$. Then the obstruction
to factoring through $\tau_{\geq t + 1}(K \otimes_A M)$ is an element
in $\Ext^t(I^nM, H_t(K \otimes_A M))$. The finite $A$-module
$H_t(K \otimes_A M)$ is annihilated by $I$. Then by
Lemma \ref{lemma-ext-annihilated}
we can after increasing $n$ assume this obstruction element is zero.
Repeating this a finite number of times we find $n$ such that
$I^nM \to K \otimes_A M$ factors through
$0 = \tau_{\geq r + 1}(K \otimes_A M)$ in $D(A)$ and we win.
\end{proof}

\end{document}
